\section{市场结构变化：从 App 中心到 Agent 接口中心}

访谈中最具争议的判断之一是“80\% app 可能消失或弱化”。更准确地说，是“许多薄交互 app 会被 agent 调度层吸收”。

\subsection{为何“App 消失”本质是交互层迁移}

当 agent 能直接操作 API 或 GUI，用户不再需要反复进入单功能 app。价值中心会从“界面拥有权”转向“能力可调用性”。

\begin{importantbox}{被重估的产品价值}
\begin{itemize}
    \item 单点 UI 价值下降；
    \item 数据可用性与 API 可靠性价值上升；
    \item 权限编排与结算机制成为新基础设施；
    \item agent-friendly 服务将获得先发优势。
\end{itemize}
\end{importantbox}

\begin{knowledgebox}{企业侧的两条转型路径}
\begin{enumerate}
    \item 把现有 app 做成稳定 API 产品；
    \item 接受 agent 作为第一入口，重构计费与风控策略。
\end{enumerate}
若两条都不做，就会被“可脚本化替代层”逐步吞噬。
\end{knowledgebox}

\begin{warningbox}{防爬策略与用户价值之间的冲突}
访谈提到平台对 bot 的强拦截。短期可防滥用，但若阻断真实用户的 agent 代理访问，长期可能把用户推向更开放的替代平台。
\end{warningbox}

\begin{figure}[H]
\centering
\includegraphics[width=0.9\textwidth]{frames/frame-023622.jpg}
\caption{产品形态讨论：Heartbeat、主动触发与“服务即接口”\protect\footnotemark}
\end{figure}
\footnotetext{视频画面时间区间：02:36:15--02:36:35。}

\subsection{本章小结}
OpenClaw 争论的实质不是“要不要 app”，而是“谁拥有任务入口”。Agent 一旦成为默认入口，软件公司的核心资产将从界面迁移到接口与数据。

